\section{Introduction}
\label{sec:intro}

Learned locomotion controllers for legged robots are trained on
distributions of terrain, and the choice of that distribution shapes what
the controller can do. Practical pipelines use hand-designed
curricula: terrain types arranged by a difficulty parameter that a robot is
promoted along once it walks far enough~\citep{rudin2022learning}, or
adaptive samplers over a parameterised terrain
family~\citep{lee2020learning}. Recent locomotion work automates this with
learning-progress~\citep{li2026lpacrl}, learned-difficulty~\citep{liu2026tacl}
and predicted-safety-cost~\citep{ovinnikov2026psc} curricula, each evaluated
with its own baselines, robot and protocol.

Unsupervised environment design (UED) studies the same problem in general
form: a teacher proposes environment parameters (\emph{levels}) to a student
so that the student generalises, typically by maximising an estimate of the
student's regret. PAIRED~\citep{dennis2020paired} learns a level generator;
PLR and its robust variant \plr~\citep{jiang2021plr,jiang2021replay} curate
randomly sampled levels by a regret score; ACCEL~\citep{parkerholder2022accel}
adds mutation of curated levels; SFL~\citep{rutherford2024noregrets} replaces
regret estimates with learnability. These methods are compared mostly on
mazes, grid worlds, car racing and the 2D BipedalWalker; libraries such as
minimax~\citep{jiang2023minimax} and JaxUED~\citep{coward2024jaxued} make
those comparisons cheap. To our knowledge no benchmark runs several UED
methods under one protocol on 3D legged locomotion, where the student is a
high-dimensional, contact-rich controller and a level is a terrain.

We build a minimal one and use it to ask whether the UED findings from 2D
domains carry over. Our contributions are:
\begin{itemize}
  \item \textbf{A two-robot benchmark} in MuJoCo MJX~\citep{todorov2012mujoco}
    with a closed-kinematic-chain quadruped (Wojtek, ${\sim}0.10$\,m
    standing height, our own physically built robot) and the open-chain
    Unitree Go1 task from MuJoCo Playground~\citep{zakka2025playground},
    scaled so that both face the same terrain relative to body size
    (\S\ref{sec:benchmark}).
  \item \textbf{An interpretable, calibrated level space}: eight terrain
    parameters plus a seed, whose ranges are set from the capability
    envelopes of pilot policies, aiming for uniform sampling that is
    neither hopeless nor trivially easy (\S\ref{sec:levelspace}).
  \item \textbf{Six methods on one student}: DR, DR-easy, \plr, ACCEL,
    ACCEL-SFL and PAIRED, all driving the same brax
    PPO~\citep{schulman2017ppo,freeman2021brax} student at matched
    environment-step budgets (\S\ref{sec:methods}).
  \item \textbf{An evaluation protocol} with cross-fitted CVaR, a
    cross-fitted solvable-subset success, per-level reliability,
    capability envelopes, a held-out structured terrain suite and a
    sim-to-sim transfer check (\S\ref{sec:protocol}), and a pilot study
    whose outcome motivated the calibration (\S\ref{sec:results},
    Appendix~\ref{app:pilot}).
\end{itemize}
A method interacts with the benchmark only through the level vectors it
proposes and whether the student trains on them, so the level space,
calibration and evaluation carry over unchanged to new methods.
\todo{Code release statement (anonymised link) and licence.}
